% =========================================================================
% secciones_es/sec06_discussion_policy.tex — Sección VI, español
% =========================================================================
\section{Discusión e implicaciones de política pública}
\label{sec:discussion}

\textbf{Resiliencia determinista para islas expuestas a huracanes.} Los sistemas insulares del Caribe no planifican contra clima estadístico sino contra tormentas con nombre. Un benchmark cuyos márgenes $N-1$ portan envolventes de bootstrap honestas en correlación, cuya mayor contingencia creíble se ejercita explícitamente y cuya resolución de estado se sitúa cómodamente dentro del ciclo de cribado EMS de 1--5~s convierte la resiliencia de una aspiración en un cómputo programado. Para un país que ha movilizado dos veces en doce meses alrededor de interrupciones nacionales \cite{ocseni}, \cite{diarioLibreHainamosa_2026}, y donde el restablecimiento tras el colapso históricamente excede de 10 a 14 horas debido a sobretensiones Ferranti ($B_{\mathrm{shunt}}$) en líneas descargadas que disparan las unidades de arranque en negro, la diferencia entre estudios heredados de horas y estados diakópticos de 2.65~ms es profunda: la partición territorial de la Ley 368-22 provee el sustrato computacional de régimen permanente para la restauración modular ascendente mediante microrredes territoriales concurrentes, mitigando las sobretensiones por efecto Ferranti asociadas a la energización de líneas en vacío y acelerando la resincronización nacional.

\textbf{La ley de planificación como frontera computacional.} La partición de este benchmark no es un artefacto de optimización: es el estatuto territorial nacional \cite{ley368}. Cuando la partición eléctrica coincide con las regiones de planificación de la Ley 368-22, cada resultado de simulación está ya expresado en las unidades en que se deciden permisos, uso de suelo y autonomía de microrredes --- y la pregunta de autonomía (¿qué territorio puede sostenerse solo tras la pérdida de una troncal?) se convierte en una propiedad del modelo y no en un estudio separado. El marco de incentivos renovables \cite{ley57} y el ciclo nacional de planificación de la expansión \cite{cne2022} adquieren, en consecuencia, un sustrato probatorio común.

\textbf{Alineación de estándares y soberanía de datos.} El benchmark habla los estándares con los que debe interoperar: intercambio de modelos consistente con CIM \cite{iec61970}, definiciones de medición de calidad de potencia de grado metrológico \cite{iec61000}, \cite{ieee1459}, práctica armónica \cite{ieee519}, práctica recomendada de métodos de estudio \cite{ieee399} e interconexión de recursos basados en inversores \cite{ieee2800}, \cite{ieee1547}. Al mismo tiempo, los archivos operativos subyacentes permanecen bajo custodia nacional: el benchmark publicado porta estructura, parámetros y arnés de validación, mientras que la capa de telemetría soberana está gobernada por instituciones nacionales. Esta separación --- ciencia abierta sobre estructura, custodia soberana sobre archivos --- es una arquitectura de política deliberada, no una omisión.

\textbf{Construcción de capacidades.} Una red nacional que puede simularse abiertamente es una red nacional que puede enseñarse abiertamente. El mismo artefacto que sirve a este artículo sirve a los laboratorios universitarios: los ejercicios doctorales y de grado se ejecutan contra el patrimonio real --- sus distancias, sus salidas, sus anclajes --- en lugar de las abstracciones de mesa de dibujo de 1962, cerrando la brecha entre el aula y la sala de control.

\textbf{Limitaciones.} Las fronteras honestas de la presente liberación son: (i) la capa de 69~kV entra mediante equivalentes estáticos, exactos en la frontera para flujo de potencia pero no para transitorios de cortocircuito dentro de las áreas condensadas; (ii) la formulación es en dominio fasorial --- los estudios de transitorios electromagnéticos requieren una extensión EMT; (iii) la telemetría llega con latencia operativa, de modo que el gemelo es causal pero no clarividente; y (iv) la liquidación de mercado queda fuera del alcance de la presente liberación, aunque la resolución horaria del horizonte es de grado de liquidación.

\textbf{Replicación para otros sistemas insulares.} Finalmente, la metodología es deliberadamente portable, y el protocolo de porteo es parte de la contribución. Como la formulación invariante no toca geografía alguna, replicar el SENI-142 para otro sistema insular requiere solo re-instanciar la cláusula territorial: (1) el patrimonio de red ensamblado bajo la misma regla nodal (una barra por nivel de tensión por subestación, transformadores como ramas); (2) el horizonte operativo ingerido bajo la misma frase de protocolo de datos con el operador del sistema sustituible nombrado; (3) la partición declarada desde el estatuto territorial rector de la jurisdicción, cualquiera que sea, con la capa jurídica agrupada en macrodominios eléctricos por topología de enlaces troncales; y (4) el ciclo de vida determinista re-sellado desde la primera compuerta. El benchmark de la República funciona así simultáneamente como artefacto y como plantilla --- un instrumento para su propio país y un contrato de reproducibilidad para el mundo insular.
